\subsection{DERIVATIONS OF AN ALGEBRA}

Let A be a graded K-algebra; for every homogeneous element \(a \in \mathbf{A}\) , let \(\operatorname{ad}_{\varepsilon}(a)\) , or simply \(\operatorname{ad}(a)\) if no confusion can arise, denote the K-linear mapping of A into A

\[
x \mapsto [ a, x ] _ {\varepsilon}
\]

(no. 4, formula (10)) which is graded of degree deg a.

PROPOSITION 8. Let A be a graded K-algebra.

\begin{enumerate}
\def\labelenumi{(\roman{enumi})}
\tightlist
\item
  For every \(\varepsilon\) -derivation \(d: \mathbf{A} \to \mathbf{A}\) and every homogeneous element \(a\) of \(\mathbf{A}\) ,
\end{enumerate}

\[
[ d, \mathrm{ad} _ {\varepsilon} (a) ] _ {\varepsilon} = \mathrm{ad} _ {\varepsilon} (d a).\tag{24}
\]

\begin{enumerate}
\def\labelenumi{(\roman{enumi})}
\setcounter{enumi}{1}
\item
  If the algebra \(\mathbf{A}\) is associative, \(\operatorname{ad}_{\varepsilon}(a)\) is an \(\varepsilon\) -derivation of \(\mathbf{A}\) of degree \(\deg(a)\) .
\item
  Suppose that \(d\) is of degree \(\delta\) , let \(\alpha = \deg a\) and write \(f = [d, \operatorname{ad}_{\varepsilon}(a)]_{\varepsilon}\) . For every homogeneous element \(x \in A\) of degree \(\xi\) , we have, by (1) (no. 1),
\end{enumerate}

\[
\begin{array}{r l} f (x) & = d (a x - \varepsilon (\alpha , \xi) x a) - \varepsilon_ {\delta , \alpha} (a (d x) - \varepsilon_ {\alpha , \delta + \xi} (d x) a) \\ & = (d a) x + \varepsilon_ {\delta , \alpha} a (d x) - \varepsilon_ {\alpha , \xi} (d x) a - \varepsilon_ {\delta + \alpha , \xi} x (d a) \\ & \quad - \varepsilon_ {\delta , \alpha} a (d x) + \varepsilon_ {\alpha , \xi} (d x) a \\ & = (d a) x - \varepsilon_ {\delta + \alpha , \xi} x (d a) = [ d a, x ] _ {\varepsilon}. \end{array}
\]

\begin{enumerate}
\def\labelenumi{(\roman{enumi})}
\setcounter{enumi}{1}
\tightlist
\item
  For all \(x\) homogeneous of degree \(\xi\) and all \(y\) homogeneous of degree \(\eta\) in \(\mathbf{A}\) ,
\end{enumerate}

\[
\begin{array}{r l} \mathrm{ad} _ {\varepsilon} (a) (x y) & = a (x y) - \varepsilon_ {\alpha , \xi + \eta} (x y) a \\ & = (a x - \varepsilon_ {\alpha , \xi} x a) y + \varepsilon_ {\alpha , \xi} x (a y - \varepsilon_ {\alpha , \eta} y a) \\ & = \mathrm{ad} _ {\varepsilon} (a) (x). y + \varepsilon_ {\alpha , \xi} x. \mathrm{ad} _ {\varepsilon} (a) (y) \end{array}
\]

taking account of (1) and the associativity of A.

When A is associative, \(\operatorname{ad}_{\varepsilon}(a)\) is called the inner \(\varepsilon\) -derivation of A defined by a.

COROLLARY. Let A be an associative graded algebra. For two homogeneous elements, a, b of A,

\[
[ \mathrm{ad} _ {\varepsilon} (a), \mathrm{ad} _ {\varepsilon} (b) ] _ {\varepsilon} = \mathrm{ad} _ {\varepsilon} ([ a, b ] _ {\varepsilon}).\tag{25}
\]

It suffices to replace \(d\) by \(\operatorname{ad}_{\varepsilon}(a)\) and \(\operatorname{ad}_{\varepsilon}(a)\) by \(\operatorname{ad}_{\varepsilon}(b)\) in (24).

If \(\deg a = \alpha\) , \(\deg b = \beta\) , formula (25) is equivalent to the following relation for every homogeneous element \(c \in \mathbf{A}\) of degree \(\gamma\)

\[
\varepsilon_ {\alpha , \gamma} [ a, [ b, c ] _ {\varepsilon} ] _ {\varepsilon} + \varepsilon_ {\beta , \alpha} [ b, [ c, a ] _ {\varepsilon} ] _ {\varepsilon} + \varepsilon_ {\gamma , \beta} [ c, [ a, b ] _ {\varepsilon} ] _ {\varepsilon} = 0\tag{26}
\]

called the Jacobi identity.

